% Part:sets-functions-relations
% Chapter: size-of-sets
% Section: schroder-bernstein

\documentclass[../../../include/open-logic-section]{subfiles}

\begin{document}

\olfileid{sfr}{siz}{sb}

\olsection{పరిమాణ భావన, ష్రోడర్--బెర్న్‌స్టైన్}

\begin{explain}
ఒక సహజమైన ఆలోచన ఇదిగో: $A$, $B$ను పరిమాణంలో మించకపోగా,
$B$ కూడా $A$ను పరిమాణంలో మించకపోతే, $A$, $B$ సమసంఖ్యాకాలు.
నిజం చెప్పాలంటే, ఈ ఆలోచన \emph{తప్పయితే}, మనం నిర్వచించిన
సమసంఖ్యాకత భావానికి సమితుల ``పరిమాణాల'' పోలికతో ఏ సంబంధమూ ఉందని
సమర్థించుకోవడం దాదాపు అసాధ్యం!{} అదృష్టవశాత్తూ, ఆ సహజమైన ఆలోచన
సరైనదే. ష్రోడర్--బెర్న్‌స్టైన్ ప్రమేయం దీన్ని సమర్థిస్తుంది.
\end{explain}

\begin{thm}[ష్రోడర్--బెర్న్‌స్టైన్]
	\ollabel{thm:schroder-bernstein}
	$\cardle{A}{B}$, $\cardle{B}{A}$ అయితే
	$\cardeq{A}{B}$.
\end{thm}

\begin{explain}
మరోలా చెప్పాలంటే, $A$ నుంచి $B$కు !!a{injection}, $B$ నుంచి $A$కు
!!a{injection} ఉంటే, $A$ నుంచి $B$కు !!a{bijection} ఉంటుంది.

అయితే ఈ ఫలితాన్ని నిరూపించడం నిజంగానే చాలా \emph{కష్టం}. నిజానికి,
కాంటర్ ఈ ఫలితాన్ని ప్రకటించినా, ఇతరులు దాన్ని నిరూపించారు.
\footnote{చరిత్ర గురించి మరింత తెలుసుకోవడానికి, ఉదాహరణకు
\citet[pp.~165--6]{Potter2004} చూడండి.}
\oliflabeldef{sfr:cardinals:card-sb:sec}{ష్రోడర్--బెర్న్‌స్టైన్‌ను
\emph{నిరూపించే} స్థితికి మనం
\olref[sfr][cardinals][card-sb]{sec}లో మాత్రమే చేరుకుంటాం.}{}%
ప్రస్తుతానికి దాన్ని మీరు నమ్మకంపై స్వీకరించవచ్చు (చేయాల్సిందే).

అదృష్టవశాత్తూ, ష్రోడర్--బెర్న్‌స్టైన్ \emph{సరైనది}; మనం నిర్వచించిన
$\cardeq{A}{B}$, $\cardle{A}{B}$ సంబంధాలకు ``పరిమాణం''తో సంబంధం
ఉందనే మన ఆలోచనను అది సమర్థిస్తుంది. అంతేకాక,
ష్రోడర్--బెర్న్‌స్టైన్ చాలా \emph{ఉపయోగకరం}. రెండు సమసంఖ్యాక
సమితుల మధ్య !!a{bijection}ను ఊహించడం కష్టంగా ఉండవచ్చు.
ష్రోడర్--బెర్న్‌స్టైన్ ప్రమేయం పోలికను రెండు సందర్భాలుగా విడగొడుతుంది;
అందువల్ల మొదటి సమితి నుంచి రెండవదానికి ఒక !!a{injection}, తిరిగి
రెండవదాని నుంచి మొదటిదానికి మరొకటి గురించి మాత్రమే ఆలోచించాలి.
\end{explain}
\end{document}
